% ==============================================================================
% Fichier  : 06_outil_aide_audit.tex
% Titre    : Conception d'un Outil d'Aide à l'Audit des SI par Décomposition
%            Hiérarchique
% ==============================================================================

\chapter{Conception d'un Outil d'Aide à l'Audit des Systèmes d'Information}
\label{chap:06_outil_aide_audit}

% ──────────────────────────────────────────────────────────────────────────────
\section{Pourquoi un outil structuré d'aide à l'audit~?}
\label{sec:pourquoi}
% ──────────────────────────────────────────────────────────────────────────────

L'audit des systèmes d'information est une démarche intellectuelle exigeante~:
l'auditeur doit simultanément maintenir une vue d'ensemble du périmètre audité,
descendre au niveau du détail opérationnel le plus fin, collecter des preuves,
formuler des risques et en anticiper les conséquences pour l'organisation.
Conduite sans support méthodique, cette démarche expose à deux écueils
symétriques et également coûteux~: l'omission de zones à risque et la
redondance d'investigations sans valeur ajoutée.

Un \emph{outil d'aide à l'audit} est un dispositif documentaire et
computationnel qui matérialise la logique de la mission d'audit sous une forme
directement opérationnelle. Il ne remplace pas le jugement professionnel de
l'auditeur~; il le libère des tâches mécaniques de mémorisation et de
structuration, pour qu'il puisse concentrer son attention sur ce qui requiert
réellement une expertise~: l'interprétation des faits et la formulation des
recommandations.

\begin{boxdef}{Outil d'aide à l'audit}
Un outil d'aide à l'audit est un référentiel structuré qui, pour un domaine
d'audit donné, décompose de manière exhaustive et ordonnée l'ensemble des
vérifications à conduire, depuis les grands axes de contrôle jusqu'aux questions
élémentaires à poser sur le terrain, en associant à chaque question le risque
qu'elle vise à détecter et la conséquence organisationnelle que l'absence de
contrôle ferait peser sur l'entité auditée.
\end{boxdef}

\subsection{Fonctions de l'outil}

Un outil d'aide à l'audit remplit quatre fonctions complémentaires.

\paragraph{Fonction de complétude.}
En déployant systématiquement la totalité des niveaux de décomposition, l'outil
garantit qu'aucun point de contrôle pertinent n'est oublié, quelle que soit
l'expérience de l'auditeur qui le met en œuvre.

\paragraph{Fonction de traçabilité.}
Chaque observation de terrain est rattachée, via la structure de l'outil, à un
point de contrôle, un objectif et un critère précis. Cette chaîne de traçabilité
est indispensable pour que le rapport d'audit soit opposable~: chaque constat
doit pouvoir être remonté jusqu'au référentiel qui le fonde.

\paragraph{Fonction de standardisation.}
Dès lors que l'outil est partagé au sein d'une équipe ou d'une institution, il
assure la comparabilité des missions d'audit conduites à des moments différents
ou par des auditeurs différents sur un même périmètre. Cette standardisation est
la condition sine qua non d'un suivi pluriannuel de la maturité organisationnelle.

\paragraph{Fonction d'automatisation.}
Implémenté dans un tableur, l'outil peut intégrer des formules de scoring, des
indicateurs de couverture, des alertes de non-conformité et des tableaux de bord
synthétiques, transformant ainsi un simple guide de questions en un véritable
instrument de pilotage de la mission.

\subsection{Ancrage normatif}

La nécessité d'un tel outil n'est pas une construction ex nihilo~: elle est
explicitement posée par les principaux référentiels internationaux de l'audit
des SI.

COBIT~2019 structure l'évaluation des processus IT en objectifs de gouvernance et
de gestion, eux-mêmes décomposés en pratiques, activités et indicateurs~; cette
arborescence est le prototype même de la décomposition hiérarchique décrite dans
ce chapitre. Les normes ISACA (CISA~Review~Manual) prescrivent une approche par
domaines et sous-domaines, avec des objectifs de contrôle définis pour chacun.
La norme ISO/IEC~27001 organise la sécurité de l'information en~114~mesures de
contrôle regroupées en~14~domaines, dont la vérification lors d'un audit suit
précisément la logique du critère, de l'élément probant et de la question
d'évaluation. Le COSO, enfin, articule le contrôle interne en cinq composantes
dont la vérification lors d'un audit suit la même progression~: assertion
(objectif de contrôle), test (critère et question), preuve (élément requis),
constat (risque), et impact (conséquence).

\begin{boxlizbiz}
\textbf{Cas Lizbiz's~: la nécessité vécue de l'outil}

Lors d'une mission d'audit de l'application de gestion des stocks de Lizbiz's,
l'équipe d'audit arrive sur site sans grille structurée. Elle identifie un
problème d'accès non autorisé aux données de production, mais ne documente pas
systématiquement les questions connexes~: la gestion des profils, la piste
d'audit, la journalisation des modifications. Le rapport conclut à une faille de
contrôle d'accès. Six mois plus tard, un audit externe révèle que les données
sensibles étaient également accessibles via un module non inventorié. La faille
était pourtant couverte par la grille de contrôle standard, elle avait été
omise faute d'un outil de décomposition systématique. L'outil d'aide à l'audit
aurait rendu cette omission impossible.
\end{boxlizbiz}

% ──────────────────────────────────────────────────────────────────────────────
\section{La méthode de décomposition hiérarchique}
\label{sec:methode}
% ──────────────────────────────────────────────────────────────────────────────

La méthode de décomposition hiérarchique est le cœur intellectuel de l'outil
d'aide à l'audit. Elle repose sur un principe simple~: tout audit, quel que soit
son type ou son périmètre, peut être entièrement représenté par une arborescence
à huit niveaux, depuis le domaine le plus général jusqu'aux conséquences les
plus concrètes de chaque défaillance de contrôle.

\begin{boxdef}{Méthode de décomposition hiérarchique}
La méthode de décomposition hiérarchique est un cadre d'analyse qui structure
tout périmètre d'audit en une succession ordonnée de niveaux d'abstraction
décroissante, chaque niveau étant entièrement dérivé du précédent selon une
relation logique de type \emph{général~$\to$~particulier}. Elle produit, pour
chaque domaine audité, une arborescence exhaustive reliant le champ d'audit
(niveau~1) aux conséquences organisationnelles des risques détectés (niveau~8).
\end{boxdef}

\subsection{Vue d'ensemble de l'arborescence}

Les huit niveaux de la décomposition sont les suivants~:

\begin{enumerate}
  \item \textbf{Domaine}~:  le champ d'audit considéré~;
  \item \textbf{Point de contrôle}~:  les grandes thématiques à examiner dans
        ce domaine~;
  \item \textbf{Objectif de contrôle}~:  ce que l'auditeur cherche à vérifier
        pour chaque point~;
  \item \textbf{Critère}~:  les conditions précises dont la réunion établit
        que l'objectif est atteint~;
  \item \textbf{Élément requis}~:  les documents ou preuves tangibles à
        collecter pour renseigner chaque critère~;
  \item \textbf{Question d'évaluation}~:  la formulation interrogative
        opérationnelle dérivée du critère~;
  \item \textbf{Risque}~:  ce que l'absence ou la défaillance du contrôle
        associé à la question fait peser sur l'entité~;
  \item \textbf{Conséquence}~:  l'impact organisationnel, opérationnel ou
        réglementaire de la réalisation du risque.
\end{enumerate}

La figure suivante représente cette arborescence~:

\begin{figure}[H]
\centering
\begin{tikzpicture}[
  font=\small,
  every node/.style={draw, rounded corners=3pt, fill=white,
                     text width=7.2cm, align=left, inner sep=5pt,
                     minimum height=0.75cm},
  arrow/.style={-Stealth, thick, color=chaptercolor},
  level/.style={sibling distance=8cm, level distance=1.1cm}
]
  \node[fill=chaptercolor!20, text=chaptercolor, font=\small\bfseries]
       {Niveau~1 \quad \textbf{Domaine}}
    child[arrow] { node[fill=enspblue!10]
       {Niveau~2 \quad \textbf{Point de contrôle}}
      child[arrow] { node[fill=enspblue!10]
         {Niveau~3 \quad \textbf{Objectif de contrôle}}
        child[arrow] { node[fill=enspgreen!10]
           {Niveau~4 \quad \textbf{Critère}}
          child[arrow] { node[fill=enspgreen!10]
             {Niveau~5 \quad \textbf{Élément requis}}
            child[arrow] { node[fill=navyblue!10]
               {Niveau~6 \quad \textbf{Question d'évaluation}}
              child[arrow] { node[fill=enspred!10]
                 {Niveau~7 \quad \textbf{Risque}}
                child[arrow] { node[fill=enspred!20]
                   {Niveau~8 \quad \textbf{Conséquence}}
                }
              }
            }
          }
        }
      }
    };
\end{tikzpicture}
\caption{Arborescence de la méthode de décomposition hiérarchique}
\label{fig:arborescence}
\end{figure}

% ──────────────────────────────────────────────────────────────────────────────
\section{Rôle et logique de chaque palier}
\label{sec:paliers}
% ──────────────────────────────────────────────────────────────────────────────

\subsection{Niveau~1. Le Domaine}

\begin{boxdef}{Domaine}
Le domaine désigne le champ d'audit considéré dans son ensemble, c'est-à-dire
l'entité, la fonction, le processus ou le système sur lequel porte la mission.
\end{boxdef}

Le domaine est le point d'entrée de toute décomposition~; il répond à la
question~: \emph{qu'est-ce qu'on audite~?} Il peut s'agir d'une application
informatique en service, de la fonction informatique (DSI), d'un projet
informatique, du support utilisateur, de la sécurité informatique, de la
fonction étude, ou de tout autre périmètre défini dans le mandat d'audit.

Le domaine ne contient en lui-même aucun contenu évaluable~; il sert
exclusivement de label d'en-tête pour la feuille de calcul correspondante dans
l'outil. Chaque domaine donne lieu à une feuille distincte.

\paragraph{Clé de conception.}
Un domaine bien délimité est un domaine dont les frontières sont acceptées par
l'audité. La lettre de mission doit formaliser cette délimitation avant le début
des travaux de terrain.

\subsection{Niveau~2. Le Point de contrôle}

\begin{boxdef}{Point de contrôle}
Un point de contrôle est une thématique majeure du domaine audité, identifiée
comme un axe de vérification cohérent et autonome. Il regroupe l'ensemble des
vérifications portant sur une même composante fonctionnelle, organisationnelle
ou technique du domaine.
\end{boxdef}

Les points de contrôle constituent le plan de l'audit~; ils répondent à la
question~: \emph{sur quels axes doit-on se prononcer~?} Ils sont définis par
référence à un référentiel reconnu~: guide d'audit des SI (IFACI/CIGREF), COBIT,
ISO~27001, guide ISACA, ou document normatif sectoriel. Leur nombre varie selon
le type d'audit~: six pour l'audit des applications en service, dix pour
l'audit de sécurité, quinze pour l'audit des projets informatiques.

\paragraph{Clé de conception.}
Un bon point de contrôle est \emph{homogène} (toutes les vérifications qu'il
recouvre concernent bien la même composante), \emph{exhaustif} (il ne laisse
pas de zone d'ombre dans sa thématique) et \emph{indépendant} des autres points
de contrôle du même domaine (pas de recouvrement).

\subsection{Niveau~3. L'Objectif de contrôle}

\begin{boxdef}{Objectif de contrôle}
Un objectif de contrôle est une assertion précise sur ce qui doit être vérifié
au sein d'un point de contrôle. Il formule, de manière affirmative et
vérifiable, l'état attendu d'une composante du système audité.
\end{boxdef}

Les objectifs de contrôle répondent à la question~: \emph{que doit-on
s'assurer~?} Ils sont formulés à l'infinitif~: \emph{S'assurer que\ldots} ou
\emph{Vérifier que\ldots}. Cette formulation n'est pas stylistique~: elle
rappelle que l'objectif engage la responsabilité de l'auditeur, il doit
apporter la preuve que la vérification a bien été conduite.

Un point de contrôle peut comporter plusieurs objectifs de contrôle. Par
exemple, le point de contrôle \emph{Organisation} dans l'audit des applications
en service comporte dix objectifs, allant de l'existence d'un comité
informatique jusqu'à la disponibilité d'un guide utilisateur.

\paragraph{Clé de conception.}
Chaque objectif doit être \emph{formulé en termes d'état} (ce qui doit exister
ou être effectif), et non en termes d'action (ce que l'auditeur va faire)~; la
formulation de l'action relève du niveau~6 (question d'évaluation).

\subsection{Niveau~4. Le Critère}

\begin{boxdef}{Critère}
Un critère est une condition élémentaire et vérifiable dont la satisfaction
contribue à établir qu'un objectif de contrôle est atteint. Il décompose
l'objectif en assertions unitaires, chacune susceptible d'être évaluée
indépendamment.
\end{boxdef}

Les critères répondent à la question~: \emph{à quelles conditions peut-on
considérer que l'objectif est rempli~?} Ils sont formulés sous la forme de
propositions assertives au présent de l'indicatif~: \emph{Le comité est présidé
par le Directeur Général~; les directions utilisatrices sont représentées~;
etc.}

La granularité du critère doit être telle qu'il ne soit ni trop large (auquel
cas il se confondrait avec l'objectif) ni trop étroit (auquel cas il deviendrait
une simple modalité). Un critère bien formulé peut être binaire~: soit il est
vérifié, soit il ne l'est pas.

\paragraph{Clé de conception.}
La somme des critères associés à un objectif doit être à la fois nécessaire et
suffisante~: si tous les critères sont satisfaits, l'objectif est atteint~; si
l'un d'eux ne l'est pas, l'objectif est partiellement ou totalement défaillant.

\subsection{Niveau~5. L'Élément requis}

\begin{boxdef}{Élément requis}
Un élément requis est un document, un enregistrement, un artefact ou une
preuve tangible dont la production ou la présentation par l'audité est nécessaire
pour que l'auditeur puisse se prononcer sur la satisfaction d'un critère.
\end{boxdef}

Les éléments requis répondent à la question~: \emph{quelles preuves faut-il
collecter~?} Ce sont les pièces du dossier d'audit~: procès-verbaux de réunion,
décisions de nomination, politiques formalisées, contrats, journaux système,
rapports d'activité, etc.

Ce niveau est le lien entre la démarche analytique de l'auditeur (les niveaux 1
à 4) et le terrain de la mission (la collecte documentaire). La liste des
éléments requis constitue, une fois compilée sur l'ensemble des critères, le
\emph{tableau de demandes de documents} à remettre à l'audité en début de
mission.

\paragraph{Clé de conception.}
Un élément requis doit être \emph{nominatif} (document spécifique, pas une
catégorie) et \emph{daté} si possible (version en vigueur, exercice concerné),
de manière à ce que l'audité sache précisément ce qui est attendu de lui.

\subsection{Niveau~6. La Question d'évaluation}

\begin{boxdef}{Question d'évaluation}
Une question d'évaluation est la formulation interrogative opérationnelle d'un
critère~; c'est la question que l'auditeur pose à l'audité, ou qu'il se pose
lui-même lors de l'examen d'un document, pour vérifier la satisfaction du
critère correspondant.
\end{boxdef}

Les questions d'évaluation répondent à la question~: \emph{comment vérifier
concrètement le critère sur le terrain~?} Elles sont formulées en propositions
interrogatives directes~: \emph{Existe-t-il un comité informatique~? Le comité
est-il présidé par le Directeur Général~?}

Ce niveau est le \emph{questionnaire d'audit}~: l'ensemble des questions
d'évaluation d'un domaine, une fois extrait de l'outil, constitue le guide
d'entretien utilisable lors des réunions avec l'audité ou lors de l'examen
documentaire.

\paragraph{Clé de conception~: la réponse trichotomique.}
Contrairement à une idée reçue, la question d'évaluation ne doit pas être
strictement binaire. L'expérience du terrain révèle qu'une troisième situation
échappe systématiquement à la dichotomie Oui/Non~: celle où l'auditeur ne
dispose pas d'une information suffisante pour se prononcer, non par manque
de diligence, mais parce que l'information elle-même n'existe pas, ne circule
pas ou n'est pas accessible au sein de l'entité auditée.

L'outil retient donc un \emph{système de réponse trichotomique}~:

\begin{itemize}
  \item \textbf{O} (\emph{Oui}), le critère est satisfait~; la preuve a été
        collectée et est probante~;
  \item \textbf{N} (\emph{Non}), le critère n'est pas satisfait~; la
        défaillance est avérée~;
  \item \textbf{PAS} (\emph{Pas Assez de Savoir}), également noté
        \textbf{NEK} (\emph{Not Enough Knowledge}), l'auditeur ne dispose pas
        d'une information suffisante pour conclure~: le document requis n'a pas
        été produit, les interlocuteurs ne peuvent pas répondre, ou les réponses
        obtenues sont contradictoires.
\end{itemize}

La modalité PAS/NEK n'est pas un aveu d'échec méthodologique~; elle est au
contraire un \emph{signal diagnostique de première importance}. Deux
interprétations organisationnelles lui sont associées~:

\begin{enumerate}
  \item \textbf{Défaillance du contrôle interne}~: si une proportion élevée de
        réponses PAS/NEK se concentre sur un point de contrôle donné, cela
        indique que les mécanismes de contrôle de ce point sont soit inexistants,
        soit non documentés, soit non appliqués, ce qui constitue en
        soi un risque, indépendamment du fait que le contrôle soit ou non en
        place.
  \item \textbf{Défaut de circulation de l'information}~: si les réponses
        PAS/NEK sont dispersées sur l'ensemble des points de contrôle et des
        interlocuteurs rencontrés, cela révèle que l'information pertinente pour
        le fonctionnement de l'organisation ne circule pas correctement~: les
        acteurs ne savent pas ce qui se fait, ce qui existe ou ce qui est décidé
        ailleurs dans leur propre organisation.
\end{enumerate}

\paragraph{Mesure de l'incertitude par l'entropie de Shannon.}
La proportion de réponses PAS/NEK dans un domaine ou sur l'ensemble d'un audit
peut être quantifiée par l'\emph{entropie informationnelle de Shannon},
appliquée à la distribution des trois modalités de réponse. Pour un domaine
comportant $n$~questions, en notant $p_O$, $p_N$ et $p_P$ les proportions
respectives de réponses O, N et PAS/NEK, l'entropie de Shannon s'écrit~:

\begin{equation}
H = -\, p_O \log_2 p_O \;-\; p_N \log_2 p_N \;-\; p_P \log_2 p_P
\quad \in [0,\; \log_2 3 \approx 1{,}585 \text{ bits}]
\end{equation}

L'entropie $H$ est nulle lorsque toutes les réponses sont identiques (certitude
complète) et maximale ($\approx 1{,}585$~bits) lorsque les trois modalités sont
équiprobables (incertitude maximale). En pratique~:

\begin{itemize}
  \item une entropie élevée sur un domaine signale une \emph{zone grise}~:
        l'auditeur ne peut pas se prononcer de manière fiable, ce qui est en
        soi un constat d'audit~;
  \item la cartographie des entropies par domaine et par point de contrôle
        constitue un \emph{indicateur de qualité de l'information} au sein de
        l'entité, complémentaire au taux de conformité.
\end{itemize}

Ce double signal, taux de conformité d'une part, entropie informationnelle
d'autre part, fait de l'outil un instrument à deux dimensions~: il mesure à
la fois \emph{ce qui ne va pas} (réponses~N) et \emph{ce qu'on ne sait pas}
(réponses~PAS/NEK), les deux dimensions pouvant pointer des pathologies
organisationnelles distinctes.

\subsection{Niveau~7. Le Risque}

\begin{boxdef}{Risque}
Dans la grille de décomposition, le risque désigne la situation de défaillance
ou d'absence de contrôle que l'auditeur identifie lorsque la réponse à une
question d'évaluation est négative. Il est formulé comme l'antithèse du critère
correspondant.
\end{boxdef}

Les risques répondent à la question~: \emph{que se passe-t-il si la réponse à
la question est négative~?} Ils sont formulés comme des assertions descriptives
de l'état de non-conformité~: \emph{Inexistence d'un comité informatique~; Le
comité n'est pas présidé par le Directeur Général~; Les directions utilisatrices
ne sont pas représentées.}

Ce niveau est \emph{mécaniquement dérivé} du niveau~6~: à chaque question dont
la réponse est Non correspond un risque formalisé. L'outil peut donc générer
automatiquement la liste des risques actifs à partir des réponses saisies par
l'auditeur, et constituer ainsi la base de la cartographie des risques de la
mission.

\paragraph{Clé de conception.}
Le risque doit être formulé au niveau de l'état de non-conformité constaté, pas
au niveau de la cause (qui relève de l'analyse de l'auditeur) ni de l'impact
(qui relève du niveau~8). La frontière avec les niveaux adjacents doit être
tenue pour que le scoring soit cohérent.

\subsection{Niveau~8. La Conséquence}

\begin{boxdef}{Conséquence}
La conséquence est l'impact organisationnel, opérationnel, financier ou
réglementaire que la réalisation du risque fait peser sur l'entité auditée.
Elle représente le \emph{préjudice potentiel} que l'absence du contrôle
concerné rend possible.
\end{boxdef}

Les conséquences répondent à la question~: \emph{pourquoi ce risque est-il
grave~?} Elles sont formulées en termes d'impacts mesurables ou plausibles~:
\emph{Incohérence des choix stratégiques~; Application non adaptée aux besoins
du métier~; Perte de données sensibles~; Non-conformité réglementaire exposant
l'entité à des sanctions.}

Ce niveau justifie, pour chaque risque identifié, la nécessité d'une
recommandation dans le rapport d'audit. La conséquence est l'argument de fond
qui permet à l'auditeur de convaincre la direction auditée de l'urgence d'une
action corrective.

\paragraph{Clé de conception.}
La conséquence doit être \emph{proportionnée}~: une conséquence surévaluée
crédite l'auditeur d'un alarmisme peu professionnel~; une conséquence
sous-évaluée banalise un constat qui mérite attention. Elle doit également être
\emph{rattachée au contexte de l'entité} et non générique~: l'absence d'un PCA
n'a pas les mêmes conséquences pour un opérateur de services critiques que pour
une PME commerciale.

% ──────────────────────────────────────────────────────────────────────────────
\section{Universalité et adaptabilité de la méthode}
\label{sec:universalite}
% ──────────────────────────────────────────────────────────────────────────────

Un trait essentiel de la méthode de décomposition hiérarchique est son
\emph{invariance structurelle}~: la succession des huit niveaux reste identique
quel que soit le type d'audit considéré. Ce qui change d'un type à l'autre,
c'est uniquement le contenu de chaque niveau, le nombre de points de
contrôle, la formulation des objectifs, la nature des éléments requis, mais
jamais la logique de décomposition elle-même.

\begin{boiteTheoreme}{Invariance structurelle de la décomposition}
Pour tout domaine d'audit $\mathcal{D}$, il existe une décomposition
hiérarchique $\mathcal{H}(\mathcal{D})$ en huit niveaux tels que~:
\begin{enumerate}
  \item Chaque élément de niveau $k$ est intégralement dérivé d'un et d'un seul
        élément de niveau $k-1$~;
  \item La réunion de tous les éléments de niveau $k$ associés à un élément
        de niveau $k-1$ couvre exhaustivement cet élément~;
  \item Deux éléments de même niveau rattachés au même parent de niveau
        $k-1$ sont disjoints.
\end{enumerate}
Cette propriété garantit que l'outil d'audit construit à partir de
$\mathcal{H}(\mathcal{D})$ est à la fois exhaustif, non redondant et
entièrement traçable.
\end{boiteTheoreme}

Cette invariance a une conséquence pratique immédiate~: \emph{apprendre la
méthode une fois suffit pour l'appliquer à n'importe quel type d'audit}. Un
auditeur qui maîtrise la décomposition sur l'audit des applications en service
peut, en changeant le contenu des cellules sans changer la structure du tableur,
construire une grille d'audit de la sécurité informatique, d'un projet ou de la
fonction DSI.

Le tableau~\ref{tab:types-niveaux} illustre la variabilité du contenu avec
l'invariance de la structure sur trois types d'audit.

\begin{table}[H]
\centering
\small
\caption{Invariance structurelle~: le même niveau pour trois domaines différents}
\label{tab:types-niveaux}
\begin{tabular}{p{2.6cm}p{3.5cm}p{3.5cm}p{3.5cm}}
\toprule
\textbf{Niveau} &
\textbf{Audit Applications} &
\textbf{Audit Sécurité} &
\textbf{Audit Projets} \\
\midrule
\textbf{1. Domaine} &
Application en service &
Sécurité informatique &
Projet informatique \\
\addlinespace
\textbf{2. Point de~contrôle} &
Organisation (6~pts) &
Politique de sécurité (10~pts) &
Planification (15~pts) \\
\addlinespace
\textbf{3. Objectif} &
S'assurer de l'existence d'un comité informatique &
S'assurer qu'une politique formalisée existe &
S'assurer de l'existence d'un planning directeur \\
\addlinespace
\textbf{4. Critère} &
Le comité est présidé par le DG &
La politique couvre l'ensemble du SI &
Le planning intègre une gestion optimale des ressources \\
\addlinespace
\textbf{5. Élément requis} &
PV de session du comité &
Document de politique de sécurité en vigueur &
Planning directeur signé et révisé \\
\addlinespace
\textbf{6. Question} &
Le comité est-il présidé par le DG~? &
La politique couvre-t-elle l'ensemble du SI~? &
Un planning directeur commun existe-t-il~? \\
\addlinespace
\textbf{7. Risque} &
Le comité n'est pas présidé par le DG &
La politique ne couvre pas l'ensemble du SI &
Absence de planning directeur commun \\
\addlinespace
\textbf{8. Conséquence} &
Décisions prises par des acteurs sans niveau décisionnel adéquat &
Zones du SI non couvertes exposées sans protection &
Incohérence dans la planification, dépassements de délais \\
\bottomrule
\end{tabular}
\end{table}

% ──────────────────────────────────────────────────────────────────────────────
\section{Structuration du fichier Excel de l'outil d'aide à l'audit}
\label{sec:excel}
% ──────────────────────────────────────────────────────────────────────────────

L'implémentation naturelle de la méthode de décomposition hiérarchique est un
classeur tableur (Microsoft Excel ou LibreOffice Calc). Sa structure reflète
directement l'arborescence à huit niveaux~: chaque feuille correspond à un
domaine (niveau~1), et chaque ligne correspond à une feuille terminale de
l'arborescence, c'est-à-dire à la combinaison complète des huit niveaux pour une
unité élémentaire de vérification.

\subsection{Architecture du classeur}

\paragraph{Organisation en feuilles.}
Le classeur comporte autant de feuilles que de domaines d'audit~: une feuille
par type d'audit (application en service, fonction informatique, projets,
support/parc, sécurité, fonction étude). Une feuille supplémentaire, nommée
\texttt{SYNTHÈSE}, agrège les indicateurs de l'ensemble des domaines.

\paragraph{Nommage des feuilles.}
Chaque feuille est nommée de manière à identifier sans ambiguïté le domaine
qu'elle couvre, selon la convention~: \texttt{[N°]\_[SIGLE\_DOMAINE]}, par
exemple~: \texttt{01\_APPLI}, \texttt{02\_DSI}, \texttt{03\_PROJETS},
\texttt{04\_SUPPORT}, \texttt{05\_SECURITE}, \texttt{06\_ETUDE}.

\subsection{Structure de chaque feuille de domaine}

Chaque feuille de domaine comporte les colonnes suivantes, dans l'ordre~:

\begin{table}[H]
\centering
\small
\caption{Structure standard d'une feuille de domaine dans l'outil d'aide à
l'audit}
\label{tab:colonnes-excel}
\begin{tabular}{clp{7.5cm}}
\toprule
\textbf{Col.} & \textbf{En-tête} & \textbf{Contenu et rôle} \\
\midrule
A & \texttt{DOMAINE} &
  Nom du domaine audité (identique sur toutes les lignes de la feuille~;
  sert à la consolidation dans la synthèse). \\
\addlinespace
B & \texttt{POINT DE CONTRÔLE} &
  Libellé du point de contrôle (niveau~2). Plusieurs lignes successives
  partagent le même libellé, regroupées visuellement par un fond coloré
  alterné. \\
\addlinespace
C & \texttt{OBJECTIF DE CONTRÔLE} &
  Libellé de l'objectif de contrôle (niveau~3), formulé à l'infinitif
  (\emph{S'assurer que\ldots}, \emph{Vérifier que\ldots}). \\
\addlinespace
D & \texttt{CRITÈRE} &
  Assertion élémentaire à vérifier (niveau~4), formulée au présent de
  l'indicatif. \\
\addlinespace
E & \texttt{ÉLÉMENT REQUIS} &
  Document ou preuve à collecter (niveau~5). \\
\addlinespace
F & \texttt{QUESTION D'ÉVALUATION} &
  Question posée sur le terrain (niveau~6), appelant une réponse
  trichotomique~: \texttt{O}~/ \texttt{N}~/ \texttt{PAS}. \\
\addlinespace
G & \texttt{RÉPONSE} &
  Cellule de saisie auditeur~: réponse \emph{trichotomique} obligatoire.
  \texttt{O}~(\emph{Oui})~: critère satisfait, preuve collectée.
  \texttt{N}~(\emph{Non})~: critère non satisfait, défaillance avérée.
  \texttt{PAS}~(\emph{Pas Assez de Savoir}, aussi noté \texttt{NEK}~: 
  \emph{Not Enough Knowledge})~: information insuffisante pour conclure,
  signalant soit une défaillance du contrôle interne, soit un défaut de
  circulation de l'information.
  C'est la \emph{seule colonne remplie manuellement} lors de la mission~;
  toutes les autres colonnes calculées en dépendent. \\
\addlinespace
H & \texttt{RISQUE} &
  Formulation du risque associé à une réponse négative (niveau~7).
  Affiché en rouge automatiquement si la réponse en col.~G est~N. \\
\addlinespace
I & \texttt{CONSÉQUENCE} &
  Impact organisationnel de la réalisation du risque (niveau~8). \\
\addlinespace
J & \texttt{CRITICITÉ} &
  Score calculé automatiquement (formule)~: pondération de la
  probabilité d'occurrence du risque et de la gravité de la conséquence.
  Exprimé sur une échelle~1--4 ou 1--9 selon le référentiel retenu. \\
\addlinespace
K & \texttt{STATUT} &
  Formule conditionnelle~: \texttt{CONFORME} si \texttt{G}~=~\texttt{O}~;
  \texttt{NON~CONFORME} si \texttt{G}~=~\texttt{N}~;
  \texttt{ZONE~GRISE} si \texttt{G}~=~\texttt{PAS}.
  Mise en forme conditionnelle~: vert~/ rouge~/ orange ambré
  respectivement. \\
\addlinespace
L & \texttt{RECOMMANDATION} &
  Cellule de saisie auditeur~: action corrective recommandée. Renseignée
  pour les statuts \texttt{NON~CONFORME} et \texttt{ZONE~GRISE}. \\
\addlinespace
M & \texttt{PRIORITÉ} &
  Urgence de la recommandation~: \texttt{IMMÉDIATE}, \texttt{COURT TERME}
  ou \texttt{MOYEN TERME}. Calculée à partir du score de criticité
  (col.~J). \\
\addlinespace
N & \texttt{COMMENTAIRE} &
  Champ libre de l'auditeur~: nuance, contexte, source de la preuve
  collectée, ou motif du PAS/NEK (document non produit, réponses
  contradictoires, interlocuteur non disponible). \\
\addlinespace
O & \texttt{ENTROPIE} &
  Calculée automatiquement au niveau du point de contrôle (ligne de
  sous-total)~: $H = -\sum_{m \in \{O,N,P\}} p_m \log_2 p_m$ où $p_m$
  est la proportion de réponses de modalité~$m$ dans le point de contrôle
  considéré. Une valeur proche de $\log_2 3 \approx 1{,}585$~bits signale
  une zone grise à fort niveau d'incertitude. \\
\bottomrule
\end{tabular}
\end{table}

\subsection{Fonctionnement des formules d'automatisation}

Trois familles de formules font de la feuille de calcul un outil actif et non
un simple formulaire statique.

\paragraph{Formule de statut (colonne~K).}
Le statut de chaque ligne est calculé automatiquement à partir de la réponse
trichotomique saisie en colonne~G~:

\begin{verbatim}
=SI(G2="O";"CONFORME";SI(G2="N";"NON CONFORME";"ZONE GRISE"))
\end{verbatim}

La mise en forme conditionnelle applique un fond vert au statut
\texttt{CONFORME}, rouge au statut \texttt{NON~CONFORME}, et orange ambré au
statut \texttt{ZONE~GRISE}. Ce dernier appel visuel signale à l'auditeur que la
ligne exige un approfondissement~: investigation complémentaire, relance de
l'audité ou recherche d'une source d'information alternative.

\paragraph{Formule de criticité (colonne~J).}
La criticité est activée pour les réponses \texttt{N} et \texttt{PAS}~: une
réponse PAS/NEK porte un risque potentiel aussi réel qu'une réponse négative
avérée, même si sa nature diffère (incertitude vs défaillance). La formule
type est~:

\begin{verbatim}
=SI(G2="O";"";SI(G2="N";PROBABILITE*GRAVITE;PROBABILITE*GRAVITE*0.7))
\end{verbatim}

Le coefficient~0,7 appliqué au score PAS/NEK traduit l'incertitude
inhérente~: le risque est réel mais non entièrement établi. Ce coefficient est
paramétrable selon la politique de risque de l'institution auditrice.

\paragraph{Formule d'entropie (colonne~O, niveau point de contrôle).}
L'entropie de Shannon est calculée sur les lignes de sous-total agrégées par
point de contrôle. En nommant \texttt{NbO}, \texttt{NbN}, \texttt{NbP} les
compteurs de réponses O, N et PAS dans le groupe, et \texttt{NbTotal} leur
somme~:

\begin{verbatim}
=SI(NbTotal=0;"";
   -SI(NbO=0;0;(NbO/NbTotal)*LOG(NbO/NbTotal;2))
   -SI(NbN=0;0;(NbN/NbTotal)*LOG(NbN/NbTotal;2))
   -SI(NbP=0;0;(NbP/NbTotal)*LOG(NbP/NbTotal;2)))
\end{verbatim}

La valeur résultante, comprise entre~0 et $\log_2 3 \approx 1{,}585$~bits,
est affichée avec une mise en forme conditionnelle à gradient~: blanc pour
$H < 0{,}5$ (certitude), jaune pour $0{,}5 \leq H < 1{,}0$ (incertitude
modérée), rouge pour $H \geq 1{,}0$ (zone grise critique).

\paragraph{Formule de priorité (colonne~M).}
La priorité de la recommandation est dérivée du score de criticité~:

\begin{verbatim}
=SI(J2>=7;"IMMÉDIATE";SI(J2>=4;"COURT TERME";"MOYEN TERME"))
\end{verbatim}

\subsection{La feuille de synthèse}

La feuille \texttt{SYNTHÈSE} agrège automatiquement, pour l'ensemble des
domaines, les indicateurs de couverture et de qualité de la mission. Elle
comporte~:

\begin{itemize}
  \item le \textbf{taux de couverture} par domaine~: proportion de questions
        auxquelles une réponse a été saisie (mesure de l'avancement de la
        mission)~;
  \item le \textbf{taux de conformité} par domaine~: proportion de réponses~O
        sur l'ensemble des questions évaluées~;
  \item le \textbf{taux d'incertitude} par domaine~: proportion de réponses
        PAS/NEK, indicateur direct de la qualité de l'information disponible
        dans l'entité auditée~;
  \item la \textbf{cartographie entropique}~: histogramme des entropies de
        Shannon par point de contrôle, permettant de visualiser d'un coup d'œil
        les zones grises les plus denses~;
  \item la \textbf{cartographie des risques}~: tableau croisé des risques
        actifs (réponses~N et PAS/NEK) par domaine et par niveau de criticité~;
  \item le \textbf{plan d'action synthétique}~: liste consolidée des
        recommandations de priorité \texttt{IMMÉDIATE} et \texttt{COURT TERME},
        extraites automatiquement des feuilles de domaine.
\end{itemize}

Ces indicateurs sont visualisés dans un tableau de bord graphique à deux
dimensions~: la dimension \emph{conformité} (ce qui ne va pas, réponses~N)
et la dimension \emph{incertitude} (ce qu'on ne sait pas, réponses~PAS/NEK),
représentées sur un radar chart à six axes (un par domaine), inséré dans la
feuille \texttt{SYNTHÈSE}.

\begin{boxlizbiz}
\textbf{Cas Lizbiz's~: trichotomie O~/ N~/ PAS appliquée au domaine Sécurité}

Lors de l'audit de sécurité de Lizbiz's, trois questions du point de contrôle
\emph{Gestion des identifiants et des mots de passe} illustrent les trois
modalités de réponse~:

\smallskip
\begin{tabular}{p{6.5cm}p{1.2cm}p{5cm}}
\textbf{Question} & \textbf{Réponse} & \textbf{Interprétation} \\
\midrule
Existe-t-il une procédure de gestion des mots de passe~? &
\texttt{N} &
Défaillance avérée~; aucun document produit. Statut~: \texttt{NON~CONFORME}. \\
\addlinespace
Les identifiants inutilisés sont-ils révoqués après un délai défini~? &
\texttt{PAS} &
L'administrateur réseau interrogé déclare qu'une procédure existe «~quelque
part~», mais ne peut en produire aucune trace. Statut~: \texttt{ZONE~GRISE}. \\
\addlinespace
Vérifie-t-on qu'il y a un seul utilisateur par identifiant~? &
\texttt{O} &
Liste des comptes actifs produite~; vérification effectuée sans anomalie. \\
\end{tabular}

\smallskip
L'entropie du point de contrôle sur ces trois questions est~:
$H = -\frac{1}{3}\log_2\frac{1}{3} - \frac{1}{3}\log_2\frac{1}{3}
        - \frac{1}{3}\log_2\frac{1}{3} \approx 1{,}585$~bits~: valeur maximale,
signalant que ce point de contrôle est une \emph{zone grise critique}. La
réponse PAS/NEK sur la révocation des identifiants n'indique pas que le
contrôle est absent~; elle indique que l'organisation ne \emph{sait pas} si
elle le pratique, ce qui est, pour un auditeur, aussi préoccupant qu'une
réponse N~: un contrôle dont personne ne peut attester l'existence est un
contrôle sur lequel on ne peut pas s'appuyer.
\end{boxlizbiz}

% ──────────────────────────────────────────────────────────────────────────────
\section{Protocole de construction de l'outil pour un nouveau domaine}
\label{sec:protocole}
% ──────────────────────────────────────────────────────────────────────────────

La méthode de décomposition hiérarchique étant universelle, la construction de
l'outil pour tout nouveau domaine d'audit suit un protocole en sept étapes.

\subsection{Étape~1~: Délimiter le domaine}

Définir précisément le périmètre de l'audit (application, fonction, projet,
processus) et s'assurer que ce périmètre est formalisé dans la lettre de
mission. Créer une nouvelle feuille dans le classeur et lui attribuer le nom
codifié du domaine.

\subsection{Étape~2~: Identifier les points de contrôle}

À partir du référentiel applicable (guide d'audit IFACI/CIGREF, COBIT, ISO,
ISACA ou référentiel sectoriel), lister les grands axes de vérification du
domaine. Renseigner la colonne~B. Vérifier que la liste est exhaustive (aucun
axe majeur oublié), homogène (chaque axe couvre bien une composante
distincte) et mutuellement exclusive (pas de recouvrement entre axes).

\subsection{Étape~3~: Décomposer chaque point en objectifs de contrôle}

Pour chaque point de contrôle, formuler les objectifs de contrôle correspondants
en utilisant systématiquement les tournures \emph{S'assurer que\ldots} ou
\emph{Vérifier que\ldots} Renseigner la colonne~C. S'assurer que chaque
objectif est un état attendu vérifiable, et non une action.

\subsection{Étape~4~: Dériver les critères}

Pour chaque objectif, décomposer l'état attendu en assertions élémentaires et
vérifiables indépendamment l'une de l'autre. Renseigner la colonne~D. Vérifier
que la somme des critères est à la fois nécessaire et suffisante à l'objectif.

\subsection{Étape~5~: Spécifier les éléments requis et les questions}

Pour chaque critère~: (i) identifier le document ou la preuve qui permettrait
d'en établir la satisfaction (colonne~E)~; (ii) reformuler le critère en
question trichotomique (O~/~N~/~PAS) posable à l'audité (colonne~F). Vérifier
que la formulation est suffisamment précise pour qu'une réponse PAS/NEK soit
distinguable d'une réponse N~: si l'audité dit «~je ne sais pas~», c'est un
PAS~; si l'audité dit «~non, ça n'existe pas~», c'est un N.

\subsection{Étape~6~: Formuler les risques et les conséquences}

Pour chaque question~: (i) formuler l'état de non-conformité correspondant à une
réponse négative (colonne~H)~; (ii) décrire l'impact organisationnel de la
réalisation de ce risque (colonne~I). Évaluer la gravité de la conséquence et
renseigner la pondération dans la plage nommée \texttt{GRAVITE}.

\subsection{Étape~7~: Saisir les formules d'automatisation et tester}

Implémenter les formules de statut (colonne~K), de criticité (colonne~J),
de priorité (colonne~M) et d'entropie de Shannon (colonne~O, niveau point de
contrôle). Appliquer les mises en forme conditionnelles~: vert/rouge/orange
ambré pour le statut, gradient blanc-jaune-rouge pour l'entropie. Relier la
feuille à la feuille \texttt{SYNTHÈSE}. Effectuer un test en saisissant des
réponses fictives couvrant les trois modalités (O, N, PAS) et vérifier que les
calculs, les colorations, les scores d'entropie et les remontées en synthèse se
comportent conformément aux spécifications.

\begin{boiteTheoreme}{Complétude de la décomposition}
La décomposition d'un domaine $\mathcal{D}$ est dite \emph{complète} si et
seulement si les trois conditions suivantes sont simultanément satisfaites~:

\begin{enumerate}
  \item \textbf{Exhaustivité}~: tout risque réel du domaine $\mathcal{D}$ est
        couvert par au moins une question d'évaluation de la grille~;
  \item \textbf{Non-redondance}~: aucune question n'est logiquement équivalente
        à une autre question de la même grille~;
  \item \textbf{Traçabilité}~: toute observation de terrain est rattachable,
        via la structure de la grille, à un point de contrôle et à un objectif
        de contrôle précis.
\end{enumerate}

Ces trois propriétés sont nécessaires et suffisantes pour garantir que le
rapport d'audit produit à partir de l'outil est exhaustif, non redondant et
opposable.
\end{boiteTheoreme}
% ------------------------------------------------------------------------------
\section{Le Tableau de Bord de Synthèse (Reporting)}
% ------------------------------------------------------------------------------

L'outil ne doit pas être une simple liste de questions. Pour LiZbiZ, il doit générer une vision macroscopique.

\begin{itemize}
    \item \textbf{Calcul du Taux de Conformité :} Moyenne pondérée des scores par domaine.
    \item \textbf{Visualisation :} Utilisation de graphiques de type "Radar" (Kiviat) pour identifier visuellement les faiblesses du SI.
\end{itemize}

\begin{boxdef}
\textbf{Indicateur de Maturité :} 
$M = \frac{\sum (Note \times Coeff)}{\sum Coeff}$. \\
Ce score permet de situer LiZbiZ sur une échelle de 0 à 5 (inspirée de l'échelle CMMI de COBIT).
\end{boxdef}
